% 第 13 题解析：光路图（TikZ 线稿）
% 光线 1 入射角 90°、折射角 60°，恰好射入摄像头
\begin{tikzpicture}[font=\small, >=Stealth]
  % 矩形介质截面
  \draw[thick] (0,0) rectangle (5,2.2);
  % 摄像头（中心小梯形）
  \fill[gray!30] (2.15,0.9) -- (2.85,0.9) -- (2.55,0.45) -- cycle;
  \draw[thick] (2.15,0.9) -- (2.85,0.9) -- (2.55,0.45) -- cycle;
  % 入射点与法线
  \coordinate (P) at (0.2,2.15);
  \coordinate (N) at (0.2,3.0);
  \coordinate (In) at (-1.5,2.55);
  \coordinate (R) at (2.5,1.6);
  \draw[dashed] (P) -- (N);
  \draw (P) -- (0.2,1.2);
  % 光线 1：入射（几乎贴表面）→ 折射（60°）→ 摄像头
  \draw[->, thick] (In) -- (P);
  \draw[->, thick] (P) -- (R);
  \node[above] at (-0.6,2.65) {光线 1};
  % 角度标注
  \pic[draw, angle radius=0.45cm, angle eccentricity=1.5, "$90^{\circ}$"]
      {angle = N--P--In};
  \pic[draw, angle radius=0.55cm, angle eccentricity=1.5, "$60^{\circ}$"]
      {angle = P--N--R};
  \node at (2.5,-0.4) {由 $n=\dfrac{\sin 90^{\circ}}{\sin 60^{\circ}}=\dfrac{2\sqrt{3}}{3}$};
\end{tikzpicture}
